\begin{helper}
Let $G$ be a finite graph, let $\Delta$ be the degree scale, let $\gamma>0$,
and let $\nu$ be an activation--priority law satisfying the fixed-activation
cube and lower-rectangle identities on an activation atom $A$.  Fix distinct
independent vertices $a,b,c\in A$, and put $Q=\{a,b,c\}$.  Let
\[
\mathcal O=\{(a,b,c),(a,c,b),(b,a,c),(b,c,a),(c,a,b),(c,b,a)\}.
\]
For an order $t=(p,q,r)\in\mathcal O$, define
\[
W_t(A)=\frac{1}{\gamma^3}\int_0^\gamma\int_z^\gamma\int_y^\gamma
\left(1-\frac{x}{\gamma}\right)^{|\{w\in A\setminus\{p,q,r\}:w\sim p\}|}
\left(1-\frac{y}{\gamma}\right)^{|\{w\in A\setminus\{p,q,r\}:w\not\sim p,\ w\sim q\}|}
\left(1-\frac{z}{\gamma}\right)^{|\{w\in A\setminus\{p,q,r\}:w\not\sim p,\ w\not\sim q,\ w\sim r\}|}
\,dx\,dy\,dz .
\]
Then
\[
\nu\{A\text{ is the activation set and every }u\in Q\text{ beats every
activated neighbor in }A\}
\le
\sum_{t\in\mathcal O}\nu\{A\text{ is the activation set}\}\,
\operatorname{ofReal}(W_t(A)).
\]
Moreover $W_t(A)\ge 0$ for every $t\in\mathcal O$.
\end{helper}
\begin{proof}
Write $F=\{\omega:\omega_1=A\}$ and $M=\nu(F)$.  For each
$t=(p,q,r)$ define the first-adjacency classes
\[
S_x=\{w\in A\setminus\{p,q,r\}:w\sim p\},\quad
S_y=\{w\in A\setminus\{p,q,r\}:w\not\sim p,\ w\sim q\},
\]
\[
S_z=\{w\in A\setminus\{p,q,r\}:w\not\sim p,\ w\not\sim q,\ w\sim r\}.
\]
They avoid the distinguished coordinates by definition.  The first class
is disjoint from the other two because their tests for adjacency to $p$
are opposite; the last two are disjoint because their tests for adjacency
to $q$ are opposite.  Every activated outside neighbor of the triple belongs
to exactly one class, by choosing its first adjacent vertex in this order.

Put $m=|S_x|$, $n=|S_y|$, and $k=|S_z|$.  The affine identity
\noderef{SamplingTripleChamberAffineIntegralIdentity} gives
\[
W_t(A)=\int_0^1\int_0^z\int_0^y x^m y^n z^k\,dx\,dy\,dz
=\frac{1}{(m+1)(m+n+2)(m+n+k+3)}>0.
\]
For completeness, the inner integral is
$y^{m+n+1}z^k/(m+1)$; integrating in $y$ gives
$z^{m+n+k+2}/((m+1)(m+n+2))$, and integrating in $z$ gives
the displayed value.  These integrations are justified by continuity of
polynomials on compact intervals, including when any exponent is zero.
All three denominators are strictly positive.  This proves, in particular,
the required nonnegativity for every order.

If $M=\infty$, the summand for $(a,b,c)\in\mathcal O$ is infinite because
$\operatorname{ofReal}(W_{(a,b,c)}(A))>0$.  The entire nonnegative sum is
therefore infinite, so the measure bound follows immediately.  Henceforth
assume $M<\infty$.

Let $B=\{\omega:\omega_2(v)\in[0,1]\text{ for all }v\in V\}$.
The lower-rectangle identity with every threshold equal to $1$ gives
$\nu(F\cap B)=M$.  Both $F$ and $B$ are measurable: activation is discrete,
and $B$ is a finite intersection of coordinate interval tests.  Since $M$
is finite, $\nu(F\setminus B)=0$.  Thus intersecting any measurable event
inside $F$ with $B$ leaves its measure unchanged.

Let $E$ be the fixed-activation survival event in the statement, and set
$H_t=E\cap\{\omega:\omega_2(p)\le\omega_2(q)\le\omega_2(r)\}$.
All these events are measurable: coordinate projections are measurable,
strict and weak real comparisons are Borel, and the survival condition
is a finite Boolean combination of such comparisons.  Every three real
priorities admit an increasing ordering among the six listed orders.
Consequently $E\subseteq\bigcup_{t\in\mathcal O}H_t$, and finite
subadditivity gives $\nu(E)\le\sum_{t\in\mathcal O}\nu(H_t)$.
No disjointness is asserted for these weak chambers.

Fix $t\in\mathcal O$ and let
\[
C_t=\{u\in[0,1]^V:u_p\le u_q\le u_r,\
u_w<u_p\ (w\in S_x),\ u_w<u_q\ (w\in S_y),\
u_w<u_r\ (w\in S_z)\}.
\]
This set is Borel by its finitely many coordinate comparisons and interval
tests.  Every point of $H_t\cap B$ projects into $C_t$: survival supplies
the strict comparison for each vertex of each first-adjacency class.
Apply \noderef{SamplingFiniteProductLowerOrthantEqualityMeasurable} with
the measurable priority projection $\omega\mapsto\omega_2$, atom $F$,
mass $M$, and Borel unit-cube subset $C_t$.  Its finite-mass and
lower-rectangle hypotheses are exactly those established or assumed here.
It yields
\[
\nu(H_t)=\nu(H_t\cap B)
\le\nu\bigl(\{\omega\in F:\omega_2\in C_t\}\cap B\bigr)
=M\operatorname{vol}(C_t).
\]
The coordinates $p,q,r$ are distinct because $t$ permutes $(a,b,c)$.
Using the disjointness and exclusions proved above,
\noderef{SamplingUnitCubeChamberVolumeIntegral} bounds
$\operatorname{vol}(C_t)$ by the $\operatorname{ofReal}$ of the ascending
simplex integral, which is $\operatorname{ofReal}(W_t(A))$ by the affine
identity already cited.  Multiplication by $M$ preserves this inequality.
Summing over the six orders and using the subadditive cover proves the
claimed bound.
\end{proof}
